\section{Related work}

\subsection{Deep Learning for Remaining Useful Life Prediction}
\label{sec:dl_rul}

Remaining useful life prediction methods can generally be grouped
into two broad categories~\cite{shang_capsule_2024, bai_towards_2024}.
The first, model-based methods, relies on mathematical modeling of
the physical degradation process and system dynamics. These approaches
are grounded in physics-based reasoning and offer accurate
threshold-based failure prediction under well-understood conditions.
However, the growing complexity of modern rotating machinery has made
it increasingly difficult to formulate accurate degradation models,
especially when the underlying mechanisms are non-linear or not fully
observable~\cite{heimes_recurrent_2008}. Moreover, model-based methods
often lack flexibility in adapting to diverse fault types and
operational conditions~\cite{shang_capsule_2024}. The second category,
data-driven methods, bypass this difficulty by learning
degradation patterns directly from sensor data, requiring less
domain-specific expertise and enabling rapid deployment without full
knowledge of the physical system~\cite{datadriven,
heimes_recurrent_2008}. Their performance, however, is strongly
conditioned on the volume, quality, and representativeness of the
available training data, a dependency that becomes critical in
industrial prognostics where run-to-failure experiments are both costly
and time-consuming~\cite{s25144481, pan2025remaining}.\\

Within the data-driven paradigm, early efforts relied on shallow
machine learning algorithms such as Artificial Neural
Networks~\cite{KHATIR2020102554}, Support Vector
Machines~\cite{TRAN2012320}, Random Forests~\cite{7508982randomforest},
and Hidden Markov Models~\cite{articleHiddenmarkov} for RUL estimation.
While these methods yielded promising initial results, they often
struggled to effectively handle the noisy, high-dimensional, and
temporally dependent nature of industrial sensor data. To address
these limitations, deep learning techniques were introduced into
the RUL prediction field. Convolutional Neural Networks (CNNs) were
among the earliest adopted, exploiting sliding-window representations
for spatial feature extraction from multivariate time series. Despite
their effectiveness in capturing local patterns, CNNs struggle to
retain long-range temporal dependencies, which are essential for
modeling the sequential nature of degradation
monitoring~\cite{shang_capsule_2024, pan2025remaining}. Recurrent
Neural Networks (RNNs) addressed this limitation through feedback
loops and internal memory states~\cite{articleRNNRUL}, enabling them
to approximate non-linear temporal dynamics. However, traditional RNNs
suffer from vanishing gradient issues that impair their ability to
learn long-term dependencies which is a key requirement in RUL
estimation~\cite{heimes_recurrent_2008}. Long Short-Term Memory
(LSTM) networks were subsequently introduced and rapidly adopted
for RUL prediction tasks, preserving information over longer time
horizons through gated mechanisms~\cite{inproceedingsLSTM}. Standard
LSTM architectures, however, tend to prioritize information from the
later parts of the input sequence, potentially ignoring critical
early-stage degradation signals that are often informative for bearing
prognostics~\cite{bai_towards_2024}.\\

To capture degradation dynamics in both temporal directions,
Bidirectional LSTM (BiLSTM) networks have emerged as a particularly
effective architecture for RUL estimation.
Ghavidel et al.~\cite{ghavidel_remaining_2026} conducted a
comparative evaluation of CNN, GRU, Transformer, and BiLSTM
architectures on bearing vibration data, demonstrating that BiLSTM
achieves the highest accuracy by capturing dependencies in both
forward and backward temporal directions. This finding is consistent
with recent bearing-specific architectures:
Wan et al.~\cite{wan_robust_2026} proposed a CNN-BiLSTM-TCN hybrid
that integrates spatial feature extraction, bidirectional temporal
modeling, and long-range dependency capture for bearing RUL on the
PRONOSTIA dataset;
Deng et al.~\cite{DENG2026112193} developed a collaborative
domain-adaptation bidirectional temporal convolution network validated
across multiple operating conditions on the XJTU-SY and PHM~2012
datasets; and
Liu et al.~\cite{liu_deep_2026} proposed a BiLSTM with self-attention
and a Bayesian linear layer for cross-domain bearing RUL prediction on
XJTU-SY. More recent work has explored
hybrid architectures that integrate recurrent and Transformer
components. Jiang et al.~\cite{s26051578} combined an extended LSTM
(xLSTM) with multi-head attention for bearing RUL prediction on
XJTU-SY, while
Yuntian et al.~\cite{YONGXING2026104313} introduced a temporal
self-attention mechanism within a dual-branch transfer network on the
same dataset. The attention mechanism
allows the model to dynamically weight the importance of individual
time steps, focusing on the most informative portions of the
degradation trajectory rather than treating all observations
uniformly. Several comprehensive reviews confirm the effectiveness
of these architectural trends: Tang et al.~\cite{tang_deep_2025}
reviewed deep learning-based hybrid models for bearing RUL
forecasting and identified attention-augmented recurrent architectures
as the current state of the art, while
Palaniappan et al.~\cite{palaniappan_systematic_2026} systematically
examined AI-based approaches across multiple bearing benchmark datasets
and reached a similar conclusion regarding the superiority of hybrid
deep learning models.\\

Although these architectures are frequently evaluated on the XJTU-SY
benchmark, their reported errors are difficult to compare directly. A
survey of recent XJTU-SY studies reveals normalised RMSE values
spanning an order of magnitude, from below 0.05 to above
0.17~\cite{QIAO2026112248, YONGXING2026104313}, a spread that cannot
reflect genuine differences in predictive capability alone. It arises
instead from unstandardised evaluation protocols: studies differ in
whether the RUL target is linear or piecewise, whether the error is
scored over the whole life or only after degradation onset, whether it
is pooled over windows or averaged per bearing, and which bearings are
held out for testing. These choices are frequently left unstated, so
that a single headline number conveys little without its accompanying
protocol. This observation motivates the explicit reporting convention
adopted in this work, in which every reported error is qualified by its
scale, scoring window, and aggregation, so that it can be interpreted
unambiguously against the literature.\\

Despite these architectural advances, a fundamental limitation
persists: the performance of all data-driven models is ultimately
bounded by the quantity and distribution of the available training
data. In bearing prognostics, run-to-failure experiments are costly,
time-consuming, and yield predominantly healthy-state observations;
the critical near-failure degradation regime, where accurate RUL
prediction matters most for maintenance safety, is systematically
underrepresented. This class imbalance has been identified as a
major open challenge across the literature.
Yang et al.~\cite{YANG2026114614} explicitly noted that insufficient
degradation data leads to overfitting and limits model generalization,
proposing a bilevel meta-task correlation network to address the
resulting few-shot learning problem on XJTU-SY.
Ji et al.~\cite{ji_semi-supervised_2026} observed that most deep
learning methods depend on supervised learning and that their accuracy
degrades substantially when labeled failure data is insufficient,
motivating a semi-supervised approach. Several reviews echo this
concern: Li et al.~\cite{s25144481} emphasized that collecting fault
data is time-consuming and costly, as some equipment cannot be allowed
to run to failure; Chao et al.~\cite{machines14040412} noted that
real-world PHM datasets are sparse, noisy, and highly imbalanced, with
limited availability of labeled failure histories; and
Jin et al.~\cite{jin_graph_2025} identified the long and
unpredictable degradation process of bearings as a fundamental barrier
to obtaining comprehensive training data, recommending data
augmentation and transfer learning as potential solutions.
Existing approaches to this problem include transfer learning across
operating conditions~\cite{DENG2026112193, YONGXING2026104313},
semi-supervised learning leveraging unlabeled
data~\cite{ji_semi-supervised_2026}, and meta-learning for few-shot
generalization~\cite{YANG2026114614}. However, none of these
strategies directly address the root cause: the absence of
sufficient near-failure training samples that faithfully represent
the degradation trajectories preceding failure. This observation
motivates the generative data augmentation approach developed in
Section~\ref{sec:gan_augmentation}, which specifically targets the
underrepresented near-failure regime through stage-conditioned
synthetic data generation.\\

A second limitation, closely related to the first, concerns the
absence of uncertainty quantification in the majority of reviewed deep
learning architectures. Most produce deterministic point estimates of
RUL with no accompanying confidence bounds.
Rathi et al.~\cite{rathi2026reliability} explicitly noted the need
for a unified framework integrating uncertainty quantification for
robust RUL prediction, and
Palaniappan et al.~\cite{palaniappan_systematic_2026} identified the
lack of probabilistic methods and interpretability as a critical gap
limiting the practical deployment of AI-based bearing prognostics.
From a maintenance decision-making perspective, a point estimate of
RUL without a calibrated prediction interval is insufficient for
risk-informed scheduling: knowing that a bearing has an estimated
RUL of 35~cycles carries fundamentally different operational
implications depending on whether the true value lies between 30 and
40 or between 10 and 60. This gap motivates the uncertainty quantification
framework presented in Section~\ref{sec:uq_rul}, which
reviews existing approaches to providing calibrated prediction
intervals for RUL estimation.


% ============================================================
%  Section 2.2 — GAN-Based Data Augmentation for RUL Prediction
% ============================================================

\subsection{GAN-Based Data Augmentation for RUL Prediction}
\label{sec:gan_augmentation}

As established in Section~\ref{sec:dl_rul}, the performance of
data-driven RUL models is fundamentally constrained by the scarcity
of near-failure training data. Generative Adversarial Networks (GANs)
have emerged as a promising solution to this problem, synthesizing
realistic degradation trajectories that augment the training set and
improve prediction accuracy. This section reviews the current
landscape of GAN-based augmentation for RUL prediction and identifies
three specific limitations of existing approaches that the proposed
framework addresses.\\

The effectiveness of GAN augmentation for RUL has been demonstrated
across diverse equipment domains and GAN architectures.
Lu et al.~\cite{9486650} proposed a jointly trained GAN-LSTM
predictor on the XJTU-SY bearing dataset, where the generator,
discriminator, and health predictor are trained simultaneously rather
than sequentially, achieving a 29\% RMSE reduction over a baseline
LSTM with statistical significance ($p = 2.26 \times 10^{-16}$).
In a follow-up study, Lu et al.~\cite{lu_joint_2022} extended this
framework with a Gaussian output layer for uncertainty estimation,
reducing the RMSE by 29.4\% and the MAE by 25\% on the same dataset.
Wang et al.~\cite{vibration7020029} combined a Variational
Autoencoder with a GAN (VAE-GAN) and employed Dynamic Time Warping
(DTW) for similarity-based weighting between generated and real
degradation trajectories, demonstrating approximately 20\%
improvement on the XJTU-SY, IMS, and pantograph datasets.
Wen et al.~\cite{wen_gru-ae-wiener_2024} proposed a Loop-GAN
architecture for bearing prognostics on both PHM~2012 and XJTU-SY,
while Cheng et al.~\cite{cheng_remaining_2023} developed a GAN with
deep autoencoder structure (GAN-DAE) for cross-condition bearing RUL
prediction on XJTU-SY. Beyond bearing applications,
Behera and Misra~\cite{behera_generative_2021} applied a Conditional
GAN with a Deep GRU predictor to the C-MAPSS turbofan dataset,
reporting at least 15\% accuracy improvement, and
Zhang et al.~\cite{9300234} proposed a Convolutional Recurrent GAN
(CR-GAN) that achieved 21--33\% error reduction across two
aero-engine operating conditions and complete error elimination for
lithium-ion battery systems. Several systematic reviews confirm the
general effectiveness of GAN augmentation for prognostics, with
reported improvements ranging from 15\% to over
50\%~\cite{11129693, khan_generative_2025, nieminen_synthetic_2026}.
It should be noted, however, that these reported gains are almost
invariably expressed as relative reductions against a study-specific
baseline rather than as absolute, protocol-qualified errors, which,
as noted in Section~\ref{sec:dl_rul}, limits their comparability
across studies.\\

Despite these encouraging results, a first limitation of the existing
literature concerns the representation of the input data fed to the
generative model. The majority of GAN-based RUL studies operate on a
\textit{scalar health indicator} (HI) rather than on raw multivariate
sensor sequences. Lu et al.~\cite{9486650, lu_joint_2022} compress
the full vibration signal into a single velocity-domain RMS value
before applying the GAN; Wang et al.~\cite{vibration7020029} extract
statistical features (mean square, variance, root mean square
amplitude) and generate scalar HI trajectories; and
Cheng et al.~\cite{cheng_remaining_2023} operate on processed
frequency-domain features rather than raw signals. This compression
discards the cross-sensor correlations that characterize multivariate
degradation dynamics---correlations between horizontal and vertical
vibration channels, between vibration and temperature, or between
sensors at different measurement points. Only a minority of studies
operate on multivariate sequences directly: Zhang et al.~\cite{9300234}
use 14 sensor channels on C-MAPSS, Xiong et al.~\cite{xiong_controlled_2023}
use 18 channels on N-CMAPSS, and
Zhang et al.~\cite{zhang_multivariate_2025} propose a dual-channel
Transformer-based conditional GAN for multivariate generation on
C-MAPSS and FEMTO. However, none of these multivariate approaches
target bearing datasets, where the sensor dimensionality is lower but
the cross-channel correlations between horizontal and vertical
vibration signals are physically significant for degradation
characterization.\\

A second limitation concerns the conditioning strategy of the
generator. Most existing GANs generate degradation trajectories
unconditionally: the generator receives random noise and produces a
sequence without explicit knowledge of which degradation regime
(healthy, early degradation, or near-failure) the output should
represent~\cite{9486650, lu_joint_2022, vibration7020029,
wen_gru-ae-wiener_2024}. This unconditional generation is
problematic because the primary motivation for augmentation is the
scarcity of \textit{near-failure} data specifically, while healthy-state
samples are already abundant. An unconditional generator may produce
sequences that predominantly resemble the majority class (healthy),
failing to address the very imbalance it was designed to correct.
A small number of studies have explored conditioning strategies:
Xiong et al.~\cite{xiong_controlled_2023} condition their
physics-informed CPI-GAN on physical degradation constraints,
Zhang et al.~\cite{9300234} incorporate degradation stage information
in their CR-GAN framework, and
Zhang et al.~\cite{zhang_multivariate_2025} condition their
Transformer-based GAN on the RUL label. Outside the GAN paradigm,
Tian et al.~\cite{11503642} recently proposed a diffusion model
conditioned on degradation stage for C-MAPSS augmentation, achieving
up to 27.3\% RMSE improvement. Nevertheless, none of these
stage-conditioned approaches have been applied to bearing datasets,
where degradation stages can be automatically derived from health
indicator trajectories without requiring external annotation.\\

A third and critical limitation is the near-complete absence of
rigorous uncertainty quantification in GAN-augmented RUL frameworks,
and more fundamentally, the absence of any study examining how the
augmentation itself affects the quality of the resulting uncertainty
estimates. Among all reviewed studies, only two incorporate any form
of uncertainty estimation.
Lu et al.~\cite{lu_joint_2022} add a Gaussian output layer that
estimates the mean and variance of the predicted health indicator,
but their own reliability curves reveal that the resulting intervals
are systematically overconfident, and no formal coverage guarantee is
provided. Zhu et al.~\cite{lubricants13040145} employ a Wiener
process with Bayesian updating for stochastic RUL estimation
alongside their CMC-GAN augmentation, but this approach relies on
parametric assumptions about the degradation dynamics that may not
hold across different equipment types.
All remaining studies produce deterministic point estimates with no
accompanying confidence
bounds~\cite{9486650, vibration7020029, 9300234,
xiong_controlled_2023, behera_generative_2021, cheng_remaining_2023,
wen_gru-ae-wiener_2024}. In every case the augmentation is evaluated
solely through its effect on point accuracy; whether enlarging the
training set with synthetic near-failure data alters the calibration
of the prediction intervals is a question that, to our knowledge, has
not been posed. This absence is consequential from a
reliability engineering perspective: a point estimate of
RUL without a calibrated prediction interval provides insufficient
information for risk-informed maintenance scheduling, where the
consequence of overestimation---predicting a component will last
longer than it actually does---is a potential safety hazard.\\

These three limitations (scalar health indicator input,
unconditional generation, and the unexamined interaction between
augmentation and uncertainty calibration) collectively constrain the
practical reliability of GAN-augmented RUL prediction. The following
sections review the uncertainty quantification landscape
(Section~\ref{sec:uq_rul}) and multi-task learning approaches
(Section~\ref{sec:multitask_rul}) that address the remaining gaps
identified in this analysis. A comprehensive comparison of all
reviewed methods is presented in Table~\ref{tab:comparison}.


\subsection{Uncertainty Quantification for RUL Prediction}
\label{sec:uq_rul}

The previous sections established that most deep learning models for RUL
prediction produce deterministic point estimates and that
GAN-based augmentation, while effective at improving point accuracy,
has been combined with uncertainty quantification in only two studies,
both using parametric methods without coverage guarantees. This section
reviews the broader landscape of uncertainty quantification (UQ)
methods for RUL prediction to situate conformal prediction within the
existing literature and to identify why its interaction with generative
augmentation remains unexplored.\\

Parametric UQ methods represent the most widely adopted approach in
the RUL literature. Bayesian Neural Networks (BNNs) have been applied
to bearing prognostics by Yang et al.~\cite{10854686} using a
GRU-BNN architecture on a custom rolling bearing dataset, by
Pan et al.~\cite{pan_two-stage_2025} using variational inference on
PRONOSTIA, by Wang et al.~\cite{wang_physics-constrained_2026} using
physics-constrained Bayesian inference on bearing data, and by
Chang et al.~\cite{chang_few-shot_2025} using Bayesian
meta-learning for few-shot bearing RUL on XJTU-SY. On turbofan
engines, Salehi et al.~\cite{salehi_bayesian_2026} developed a
Bayesian Informer achieving 93\% coverage, and
Zeng et al.~\cite{zeng_2026} proposed an attention-based Bayesian
ensemble reporting a PICP of 99.86\% on C-MAPSS. Monte Carlo Dropout
(MC Dropout) is the second most common parametric approach:
Ge et al.~\cite{ge_bearing_2025} combined MC Dropout with a Gaussian
output layer for bearing RUL on XJTU-SY and PHM~2012, reporting PICP
and PINAW metrics;
Ye et al.~\cite{11131520} applied MC Dropout with contrastive
learning on XJTU-SY; and
Yilma et al.~\cite{yilma_deep_2026} combined evidential deep
learning with MC Dropout on C-MAPSS, reporting both PICP and PINAW.
Gaussian output layers, which estimate the mean and variance of the
predicted distribution directly, have been used by
Zhang et al.~\cite{11398369} and
Chen et al.~\cite{chen_probabilistic_2026}, with the latter
reporting both PICP and MPIW on C-MAPSS. A common limitation across
all parametric methods is their reliance on distributional
assumptions, typically Gaussian, that may not hold for the
non-linear, non-stationary degradation dynamics characteristic of
industrial bearings. When these assumptions are violated, the
resulting prediction intervals can be miscalibrated:
Vashishtha et al.~\cite{s26072230} reported a PICP of only 0.594
on the IMS bearing dataset using a VAE-based probabilistic model,
indicating severe under-coverage, and
Zhou et al.~\cite{zhou_self-attention_2026} explicitly acknowledged
that MC Dropout underestimates the true variability of degradation
data, requiring post-processing corrections to mitigate
overconfidence. Similarly, Lu et al.~\cite{lu_joint_2022}, the only
study combining GAN augmentation with parametric UQ, showed through
reliability curves that their Gaussian layer produces systematically
overconfident intervals on the XJTU-SY bearing dataset.\\

Quantile regression and VAE-based approaches offer partially
distribution-free alternatives but introduce their own limitations.
Zhang et al.~\cite{zhang_explainable_2025} applied a gated quantile
regression network with kernel density estimation to XJTU-SY and IMS
bearing data, reporting an average coverage of roughly 94\% across the
two datasets---one of the few bearing UQ studies with rigorous interval
evaluation. While quantile regression does not assume a
specific output distribution, it requires careful selection of
quantile levels and does not provide a finite-sample coverage
guarantee in the conformal sense. VAE-based methods, including the
Dynamical VAE of Star and McKee~\cite{star_dynamical_2025} and the
Gaussian Process VAE of Ding et al.~\cite{10738265}, generate
uncertainty estimates through latent space sampling. However, these
approaches are architecture-specific, in that the UQ mechanism is
inseparable from the generative model, making it impossible to
independently evaluate the contribution of data augmentation versus
uncertainty estimation---a methodological limitation that the proposed
framework avoids by deliberately separating the generative component
(the GAN) from the UQ component (conformal prediction).\\

Conformal prediction has recently emerged as a principled alternative
that provides distribution-free, finite-sample coverage guarantees
regardless of the underlying model architecture.
Javanmardi and H{\"u}llermeier~\cite{javanmardi_conformal_2023}
conducted the most comprehensive study to date, comparing Split
Conformal Prediction (SCP), Conformalized Quantile Regression (CQR),
and non-exchangeable SCP on the C-MAPSS dataset with deep
convolutional neural networks and gradient boosting as base
predictors. Their key finding, that conformal interval quality depends
directly on the nonconformity scores of the underlying regression
model, provides the theoretical basis for the interaction examined in
this paper: altering the base predictor's error distribution through
GAN augmentation should in turn alter the quality of its conformal
intervals, whether by narrowing them at fixed coverage or by bringing
their empirical coverage closer to the nominal target. Subsequent
conformal prediction studies
have expanded the methodological repertoire:
Xu and Zio~\cite{10381131} demonstrated model-agnostic split CP
across six different base predictors (SVR, RF, GB, CNN, DNN, LSTM)
on C-MAPSS at 90\% coverage;
Piao et al.~\cite{10816036} proposed CRULP, combining
uncertainty-aware CQR with ensemble cross-conformal prediction;
Piao et al.~\cite{11163867} addressed covariate shift with an
adaptive conformalized predictor;
Robinson et al.~\cite{robinson_remaining_2026} integrated risk-aware
conformal prediction with maintenance cost optimization;
Moccardi et al.~\cite{fi17060244} compared naive, bootstrapping, and
weighted conformal prediction for IoT-based predictive maintenance;
and Hattar and Vincent~\cite{11044907} applied conformal prediction
with Random Forest and LightGBM ensembles. Recent hybrid
approaches have combined conformal prediction with other UQ methods:
Alomari et al.~\cite{11488562} integrated MC Dropout with conformal
prediction in a multi-task framework on N-CMAPSS, achieving 97.3\%
empirical coverage;
Yang et al.~\cite{yang_uncertainty_2026} combined MC Dropout with
conformal prediction using a Transformer on C-MAPSS at 95\% coverage;
Xu et al.~\cite{xu_novel_2026} paired Bayesian deep learning with
split conformal prediction on C-MAPSS and a lithium-ion battery
dataset; and
Wen et al.~\cite{en19020414} combined MC Dropout with inductive
conformal prediction for battery RUL.\\

Only very recently has conformal prediction been applied to bearing
RUL specifically, in two concurrent 2026 studies.
Bott et al.~\cite{11373138} combined a physics-based degradation
simulator with conditional normalizing flows for sim-to-real feature
alignment, training tabular predictors exclusively on synthetic data
and wrapping their output with split conformal prediction on the
XJTU-SY dataset; their symmetric conformal intervals attain an
empirical coverage of 0.919 against a 90\% nominal target, with a mean
width of 0.654.
Wang et al.~\cite{wang_uncertainty-aware_2026} proposed a
physics-informed framework in which a variational autoencoder
constructs a monotonic health index, a Transformer predictor is
trained with a monotonicity constraint, and Monte Carlo dropout
intervals are recalibrated post hoc through conformal prediction;
on XJTU-SY their coverage rises from 0.484 before calibration to 0.945
after, confirming that the raw intervals were severely overconfident
and that conformal calibration corrects this. These two studies
establish that conformal calibration is both applicable and necessary
for bearing RUL, and the present work builds directly on that footing.
Two features nonetheless distinguish them from the approach proposed
here. First, both report a single aggregate coverage over the whole
test population and do not examine whether coverage holds within the
near-failure regime specifically, where it matters most for maintenance
safety and where, as shown in Section~\ref{sec:windowing}, the
heterogeneity of residuals across degradation stages places the
strongest demands on the calibration procedure. Second, and more
fundamentally, both calibrate a model trained on real data only (Wang
et al.) or on physics-simulated data (Bott et al.); neither employs
generative adversarial augmentation, and neither investigates how such
augmentation affects conformal coverage.\\

Three gaps therefore persist in the conformal prediction literature for
RUL estimation. First, no study examines the effect of generative
adversarial augmentation on conformal coverage: existing work
conformalizes models trained on real or physics-simulated data, leaving
unexplored whether enriching the training set with synthetic
near-failure data alters interval calibration, despite the theoretical
argument of Javanmardi and
H{\"u}llermeier~\cite{javanmardi_conformal_2023} that interval quality
is governed by the base predictor's nonconformity scores. Second, the
two existing bearing conformal studies report only aggregate coverage,
without the stage-conditional calibration needed to guarantee coverage
within the near-failure regime whose residual distribution differs
sharply from the healthy one. Third, no study combines conformal
prediction with multi-task learning on bearing data, which would allow
the conformal layer to benefit from a representation shared with an
auxiliary degradation-stage task.


\subsection{Multi-task Learning for RUL Prediction}
\label{sec:multitask_rul}

Beyond uncertainty quantification, another avenue for improving
RUL prediction accuracy lies in the learning architecture itself.
Multi-task learning (MTL) offers a principled mechanism for
extracting richer representations by jointly training a shared
encoder on multiple related tasks, enabling the model to learn
more discriminative features than either task alone would produce.
In the context of bearing prognostics, the most natural auxiliary
tasks are fault mode classification and degradation stage
identification, both of which share informative features with RUL
regression: the temporal patterns that distinguish fault types or
degradation stages are often the same patterns that characterize
degradation severity. This section reviews the growing body of
multi-task learning research for RUL prediction.\\

The effectiveness of multi-task learning for RUL has been
demonstrated through two dominant formulations. The first combines
RUL regression with fault mode classification.
Liu et al.~\cite{8709822} proposed the JL-CNN, an early
multi-task architecture with shared convolutional layers and a joint
loss function for simultaneous fault recognition and RUL prediction
on IMS bearing data, achieving a 24.9\% RMSE reduction over a
single-task CNN baseline.
Xu et al.~\cite{xu_novel_2025} extended this paradigm with a fault
mode feature separation network, reporting an 11.87\% RMSE
improvement through shared representation learning on both C-MAPSS
and IMS bearing datasets.
Zhang et al.~\cite{10216122} proposed a parallel feature fusion
network for simultaneous bearing fault detection and RUL prediction
on IMS, achieving 32.6\% RMSE improvement over a CNN-LSTM baseline
without multi-task learning.
More recently, Lai et al.~\cite{LAI2026112091} developed FR-LLM, a
unified multi-task large language model for joint fault diagnosis
and RUL prediction on the XJTU-SY and IMS bearing datasets,
demonstrating that multi-task learning consistently outperforms
single-task approaches in accuracy and generalization.
Liu et al.~\cite{11434241} proposed MSDA-Net, combining a
sensor-wise CNN with a BiLSTM and dual attention for joint failure
classification and RUL estimation on C-MAPSS, explicitly
identifying the absence of uncertainty quantification as a
limitation and recommending Bayesian methods for future work.\\

The second formulation combines RUL regression with health state or
degradation stage assessment. This approach has gained particular
traction in the reliability engineering community.
Zhang et al.~\cite{zhang_health_2022} proposed a BiGRU with
multi-gate mixture-of-experts (MMoE) for joint health status
assessment and RUL prediction of aero-engines, using homoscedastic
uncertainty for task weighting.
Zhang et al.~\cite{zhang_multi-task_2023} extended this with
MTLTrans, a Transformer-based architecture that jointly optimizes
three tasks---RUL prediction, state-of-health assessment, and
performance degradation prediction, achieving an 18.83\% RMSE
improvement across 12 cross-domain scenarios on C-MAPSS.
Zhou and Wang~\cite{zhou_adaptive_2024} developed an adaptive
multi-scale feature fusion model with mixture-of-experts for joint
health assessment and RUL prediction, also employing homoscedastic
uncertainty weighting for multi-task optimization.
Gao et al.~\cite{gao_bmtr-former_2026} proposed BMTR-former, a
Bayesian multi-task Transformer that combines RUL prediction with
signal reconstruction as a self-supervised auxiliary task,
specifically targeting data scarcity through sliding window
augmentation and Bayesian inference for uncertainty estimation.
On bearing data specifically,
Wang and Wu~\cite{10295667} constructed a multi-task framework for
joint health stage identification and RUL prediction on XJTU-SY,
using fuzzy clustering for automatic stage label derivation.
Chen et al.~\cite{11363290} proposed a conflict-free multi-task
framework for joint health assessment and RUL prediction on both
XJTU-SY and PHM~2012, using wavelet-based first prediction time
detection for automatic label construction. Most recently,
Zeng et al.~\cite{11451867} developed a convolutional recurrent
residual network with attention for joint degradation stage
identification and continuous RUL regression on a large-scale
bearing dataset, with a weighted MSE loss that emphasizes the
near-failure regime, an approach that addresses the same
data imbalance problem as the proposed GAN augmentation, but
through loss weighting rather than synthetic data generation.\\

Despite these advances, three critical gaps persist in the
multi-task RUL literature. The first concerns the combination of
multi-task learning with generative data augmentation. Among all
reviewed studies, only two incorporate any form of generative
augmentation.
Behera et al.~\cite{10107629} proposed an LS-R-ACGAN multi-task
framework combining GAN-based augmentation with joint RUL and
health assessment prediction on C-MAPSS, but reported inflated
improvement metrics (354\%) against unspecified baselines, used no
uncertainty quantification, and validated exclusively on turbofan
engine data, not bearing datasets.
Guan et al.~\cite{guan_enhanced_2026} combined multi-task learning
with VAE-based augmentation for semi-supervised RUL prediction on
C-MAPSS, but again without bearing validation or uncertainty
quantification. All remaining studies train exclusively on real data,
leaving the multi-task encoder vulnerable to the same near-failure
data scarcity that limits single-task models. The proposed framework
addresses this gap by training the multi-task BiLSTM on
GAN-augmented data conditioned on degradation stage, simultaneously
enriching the training set for both the RUL regression and stage
classification heads.\\

The second gap concerns uncertainty quantification. Among the
reviewed multi-task studies, only a minority incorporate any form
of UQ: Wang et al.~\cite{wang_multi-task_2024} use mixture density
networks on IMS bearing data;
Gao et al.~\cite{gao_bmtr-former_2026} employ Bayesian deep learning
on C-MAPSS; and
Zhang et al.~\cite{zhang_health_2022} and
Zhou and Wang~\cite{zhou_adaptive_2024} use homoscedastic uncertainty
weighting for task balancing rather than for producing prediction
intervals. All of these are parametric approaches that assume
specific output distributions. No multi-task RUL study applies
conformal prediction to provide distribution-free,
coverage-guaranteed prediction intervals, with one exception:
Alomari et al.~\cite{11488562} recently combined multi-task learning
with conformal prediction on N-CMAPSS, achieving 97.3\% empirical
coverage, but without generative augmentation and only on turbofan
engine data. The proposed framework extends this direction by
combining conformal prediction with both multi-task learning and GAN
augmentation on bearing datasets.\\

The third gap concerns the choice and derivation of auxiliary task
labels. The majority of fault-based multi-task studies require
manually annotated fault mode
labels~\cite{8709822, LAI2026112091, 11434241}, which are expensive
to obtain in industrial settings and unavailable for many bearing
datasets. Health state formulations avoid this problem through
automatic label derivation, using fuzzy
clustering~\cite{10295667}, wavelet-based detection~\cite{11363290},
or threshold-based
segmentation~\cite{zhang_health_2022, 11451867}, but no existing
study derives stage labels directly from the health indicator
trajectory using physically meaningful thresholds (first prediction
time and end-of-life) as proposed in this work. This approach
produces degradation stage labels that are automatically constructed,
physically interpretable, and directly aligned with the
stage-conditioned GAN augmentation of
Section~\ref{sec:gan_augmentation}: the GAN generates synthetic
samples targeting the underrepresented near-failure stage, while the
multi-task classifier learns to recognize that same stage, creating a
synergistic loop between augmentation and multi-task learning.


\subsection{Summary and Research Gaps}
\label{sec:summary_gaps}

The preceding review reveals that while substantial progress has been
made in each individual research direction, no existing study
simultaneously addresses the compounding limitations that
constrain the reliability of bearing RUL prediction in
safety-critical industrial settings.\\

Section~\ref{sec:dl_rul} established that deep learning architectures,
particularly BiLSTM networks with attention mechanisms, have
become the state of the art for RUL estimation on bearing datasets
such as XJTU-SY. However, all reviewed models are bounded by the
quantity and distribution of their training data: the critical
near-failure regime is systematically underrepresented in
run-to-failure experiments, and the reported errors are difficult to
compare across studies because evaluation protocols are rarely stated
in full.\\

Section~\ref{sec:gan_augmentation} showed that GAN-based data
augmentation effectively mitigates near-failure data scarcity, with
reported improvements of 15--30\% in RMSE. Yet existing GAN
approaches operate predominantly on scalar health indicators rather
than multivariate sensor sequences, generate degradation trajectories
without conditioning on degradation stage, and evaluate the
augmentation exclusively through point accuracy, leaving its effect on
the calibration of prediction intervals entirely unexamined.\\

Section~\ref{sec:uq_rul} demonstrated that uncertainty quantification
for RUL prediction is a rapidly growing field, with parametric methods
(Bayesian networks, MC Dropout, Gaussian layers) widely adopted but
limited by distributional assumptions that lead to miscalibration.
Conformal prediction offers a distribution-free alternative with
finite-sample coverage guarantees. It has only very recently been
applied to bearing RUL, by two concurrent 2026 studies, and in both
cases to a model trained without generative augmentation, with coverage
reported only in aggregate rather than within the critical near-failure
regime.\\

Section~\ref{sec:multitask_rul} established that multi-task learning
improves RUL accuracy through shared representation learning,
with degradation stage classification emerging as a particularly
effective auxiliary task for bearing applications. However, no
existing multi-task framework simultaneously incorporates generative
augmentation to address near-failure data scarcity and conformal
prediction to provide coverage-guaranteed intervals.\\

Table~\ref{tab:comparison} provides a systematic comparison of the
proposed framework against the most relevant existing methods across
seven dimensions. The proposed framework is, to our knowledge, the
first to integrate all three components in a single pipeline for
bearing RUL: (1)~a conditional GAN generates synthetic multivariate
sensor sequences conditioned on degradation stage, specifically
targeting the underrepresented near-failure regime; (2)~a multi-task
BiLSTM with attention jointly predicts RUL and degradation stage,
where the auxiliary classification task forces the shared encoder to
learn stage-discriminative features that improve accuracy in the
critical late-degradation period; and (3)~stage-conditional conformal
prediction wraps the RUL regression output to produce distribution-free,
coverage-guaranteed prediction intervals evaluated per degradation
stage through PICP and MPIW. Beyond this integration, the central
empirical question examined in this work has not been posed in the
prior literature: whether stage-conditioned generative augmentation
improves the quality of the conformal intervals---and in particular the
reliability of coverage in the near-failure regime---over and above its
effect on point accuracy.


% ============================================================
%  Qualitative comparison table (seven dimensions)
% ============================================================

\begin{table}[!htbp]
\centering
\caption{Comparison of the proposed framework against existing methods.}
\label{tab:comparison}
\footnotesize
\renewcommand{\arraystretch}{1.15}
\setlength{\tabcolsep}{2.5pt}
\begin{tabular}{lc ccccccc}
\toprule
\textbf{Study} & \textbf{Year}
  & \rotatebox{70}{\textbf{GAN augment.}}
  & \rotatebox{70}{\textbf{MV input}}
  & \rotatebox{70}{\textbf{Stage-cond.}}
  & \rotatebox{70}{\textbf{Multi-task}}
  & \rotatebox{70}{\textbf{UQ method}}
  & \rotatebox{70}{\textbf{Coverage guar.}}
  & \rotatebox{70}{\textbf{Bearing data}} \\
\midrule

% --- GAN augmentation ---
Lu et al.~\cite{9486650}                    & 2021 & \checkmark & \texttimes & \texttimes & \texttimes & --   & \texttimes & \checkmark \\
Lu et al.~\cite{lu_joint_2022}              & 2022 & \checkmark & \texttimes & \texttimes & \texttimes & GL   & \texttimes & \checkmark \\
Zhang et al.~\cite{9300234}                 & 2021 & \checkmark & \checkmark & \checkmark & \texttimes & --   & \texttimes & \texttimes \\
Xiong et al.~\cite{xiong_controlled_2023}   & 2023 & \checkmark & \checkmark & \checkmark & \texttimes & --   & \texttimes & \texttimes \\
Zhu et al.~\cite{lubricants13040145}        & 2025 & \checkmark & \checkmark & \texttimes & \texttimes & WP   & \texttimes & \texttimes \\
\midrule

% --- Conformal prediction ---
Javanmardi~\cite{javanmardi_conformal_2023} & 2022 & \texttimes & \checkmark & \texttimes & \texttimes & CP   & \checkmark & \texttimes \\
Piao et al.~\cite{10816036}                 & 2025 & \texttimes & \checkmark & \texttimes & \texttimes & CP   & \checkmark & \texttimes \\
Alomari et al.~\cite{11488562}              & 2026 & \texttimes & \checkmark & \texttimes & \checkmark & CP   & \checkmark & \texttimes \\
Wang et al.~\cite{wang_uncertainty-aware_2026} & 2026 & \texttimes & \checkmark & \texttimes & \texttimes & CP & \checkmark & \checkmark \\
Bott et al.~\cite{11373138}                 & 2026 & \texttimes & \checkmark & \texttimes & \texttimes & CP   & \checkmark & \checkmark \\
\midrule

% --- Multi-task learning ---
Liu et al.~\cite{8709822}                   & 2020 & \texttimes & \checkmark & \texttimes & \checkmark & --   & \texttimes & \checkmark \\
Wang et al.~\cite{wang_multi-task_2024}     & 2024 & \texttimes & \checkmark & \texttimes & \checkmark & MDN  & \texttimes & \checkmark \\
Zeng et al.~\cite{11451867}                 & 2026 & \texttimes & \checkmark & \texttimes & \checkmark & --   & \texttimes & \checkmark \\
Behera et al.~\cite{10107629}               & 2024 & \checkmark & \checkmark & \texttimes & \checkmark & --   & \texttimes & \texttimes \\
Gao et al.~\cite{gao_bmtr-former_2026}      & 2026 & \texttimes & \checkmark & \texttimes & \checkmark & BDL  & \texttimes & \texttimes \\
\midrule

\rowcolor[RGB]{198,224,180}
\textbf{Proposed}                            & \textbf{2026}
  & \checkmark & \checkmark & \checkmark & \checkmark
  & \textbf{CP} & \checkmark & \checkmark \\

\bottomrule
\end{tabular}

\vspace{4pt}
\raggedright\scriptsize
GL: Gaussian layer. WP: Wiener process. CP: conformal prediction.
MDN: mixture density network. BDL: Bayesian deep learning.
--: no uncertainty quantification.
\end{table}
